% 20260513_Assingment_ComputationalMathematics_A
\documentclass[dvipdfmx]{jsreport}
\usepackage{soupreamble}

\title{計算数学A 第1回レポート課題}
\author{M1 6012 川村聡一郎}
\date{}

\begin{document}
\maketitle

\begin{tcolorbox}
	$\fbox{1}$ $f(n)= O(g(n))$かつ$g(n)= O(h(n))$ならば
\begin{equation*}
	f(n)= O(h(n))
\end{equation*}
\end{tcolorbox}

\begin{souanswer} %答案はここに書く
	$\fbox{1}$ $f(n)\le Ch(n)$を示す．\\
	$\exists c_1, c_2\in \R_{\ge 0},\ \exists n_0, n_1\in \Z_{\ge 0},\ n\ge \max\{n_0, n_1\}$に対して
\begin{equation}
	f(n)\stackrel{f(n)= O(g(n))}{\le} c_1g(n)\stackrel{g(n)= O(h(n))}{\le} c_1c_2h(n)
\end{equation}
	$C= c_1c_2$とおくと，$f(n)\le Ch(n)$．よってbig- $O$の定義より$f(n)= O(h(n))$が成り立つ．
\end{souanswer}


\begin{tcolorbox}
	$\fbox{2}$ 以下が同値であることを示せ．
\begin{enumerate_alpha}
	\item $f(n)= \Theta(g(n))$
	\item $f(n)= O(g(n))$かつ$f(n)= \Omega(g(n))$
\end{enumerate_alpha}
\end{tcolorbox}

\begin{souanswer} %答案はここに書く
\begin{enumerate}
	\item[$\Rightarrow$] 
	仮定よりある正の定数$c_1, c_2> 0$と整数$n_0\ge 0$が存在して，すべての$n\ge n_0$に対して
\begin{equation}
	c_1g(n)\le f(n)\le c_2g(n)
\end{equation}
	が成り立つ．この不等式を上からの評価「$f(n)\le c_2g(n)$」と下からの評価「$c_1g(n)\le f(n)$」に分ける．
	すると上からの評価は$f(n)= O(g(n))$，下からの評価は$f(n)= \Omega(g(n))$の定義に一致する．

	\item[$\Leftarrow$] 
	$O(g(n)), \Omega(g(n))$それぞれの定義より
\begin{align}
	O(g(n))&= \{h(n); \exists c\in \R_{\ge 0}, \exists n_0\in \Z_{\ge 0}, n\ge n_0\Longrightarrow h(n)\le cg(n)\}\\
	\Omega(g(n))&= \{h(n); \exists c'\in \R_{\ge 0}, \exists n_0\in \Z_{\ge 0}, n\ge n_0\Longrightarrow c'g(n)\le h(n)\}
\end{align}
	これらを組み合わせると
\begin{align}
	O(g(n))\cap \Omega(g(n))&= \{h(n); \exists c, c'\in \R_{\ge 0}, \exists n_0\in \Z_{\ge 0}, n\ge n_0\Longrightarrow c'g(n)\le h(n)\le cg(n)\}\\
	&= \Theta(g(n))
\end{align}
	以上より「$f(n)= O(g(n))$かつ$f(n)= \Omega(g(n))$」と「$f(n)= \Theta(g(n))$」は同値．
\end{enumerate}
\end{souanswer}

\begin{tcolorbox}
	$\fbox{3}$ 次を示せ．
\begin{enumerate}
	\item $n^2\cdot 2^n\ne \tilde{O}_1(2^n)$
	\item $n^2\cdot 2^n= \tilde{O}_2(2^n)$
\end{enumerate}
\end{tcolorbox}

\begin{souanswer} %答案はここに書く
\begin{enumerate}
	\item $s> 0$とする．$n^s\cdot 2^n= \tilde{O}_1(2^n)$と仮定すると，$\tilde{O}_1$の定義より正の定数$c, n$と非負定数$k$が存在し，すべての$n\ge n_0$に対して，
\begin{equation*}
	n^s\cdot 2^n\le c\cdot 2^n\cdot (\log{n})^k
\end{equation*}
	が成立する．両辺を$2^n$で割ると
\begin{equation*}
	n^s \le c\cdot (\log{n})^k
\end{equation*}
	さらに両辺を$(\log{n})^k$で割ると，
\begin{equation*}
	\frac{n^s}{(\log{n})^k}\le c
\end{equation*}
	ここで$n\to \infty$の極限を考えると，いずれも$n\to \infty$で発散する．
	これは仮定に矛盾するため，$n^2\cdot 2^n\ne \tilde{O}_1(2^n)$が成り立つ．

	\item ある整数$k$に対して
\begin{equation}
	n^2\cdot 2^n= O(2^n(\log{2^n})^k)
\end{equation}
	が成り立つことを示せばよい．すなわち正の定数$c> 0$，非負整数$n_0$が存在し，すべての$n\ge n_0$に対して，
\begin{equation}
	n^2\cdot 2^n\le c\cdot 2^n(\log{2^n})^k
\end{equation}
	が成り立つことを示せばよい．不等式の右辺は$\log{2^n}= n\log{2}$と書き換えることができ，
\begin{equation}
	n^2\cdot 2^n\le c\cdot n^k(\log{2})^k\cdot 2^n
\end{equation}
	すべての$n$に対して$2^n> 0$だから両辺を$2^n$で割ると，
\begin{equation}
	1\le c\cdot n^k(\log{2})^k
\end{equation}
	ここで$k= 2$とすると，
\begin{equation}
	1\le c\cdot n^2(\log{2})^2
\end{equation}
	また$n\ge 1$とすると$n^2> 0$となるから，
\begin{equation}
	1\le c(\log{2})^2
\end{equation}
	$c= \frac{1}{(\log{2})^2}$ととればこの不等式を満たす．\\
	以上より$k= 2, c= \frac{1}{(\log{2})^2}, n_0= 1$とすると，すべての$n\ge n_0$に対して，
\begin{equation}
	n^2\cdot 2^n\le c\cdot 2^n(\log{2^n})^k
\end{equation}
	であるから，$n^2\cdot 2^n= \tilde{O_2}(2^n)$が成り立つ．
\end{enumerate}
\end{souanswer}



\begin{tcolorbox}
	$\fbox{4}$ $n!= o(n^n)$を示せ．
\end{tcolorbox}

\begin{souanswer} %答案はここに書く
\begin{align*}
	\frac{n!}{n^n}
	&= \frac{n\cdot (n- 1)\cdot (n- 2)\cdots 1}{n^n}\\
	&\le \frac{n\cdot n\cdot n\cdots 1}{n^n}\\
	&\le \frac{1}{n}\\
	&= 0\ (n\to \infty)
\end{align*}
\end{souanswer}
\end{document}
